% Generated by scripts/assemble.py; corrections belong in source/corrections.json.
\examhead{7}{北京交通大学 2026 年数学分析试卷(回忆版)}{13}
% p014-b20; source PDF page 14
\blockanchor{p014-b20}
(注.题目来自考生回忆,部分表述或数据可能有误)\par
% p014-b21; source PDF page 14
\par\Needspace{6\baselineskip}
\blockanchor{p014-b21}
\question{1}已知 $a_1, b_1$ 为正实数 $(a_1 > b_1)$, 且满足\par
\nopagebreak[4]
% p014-b22; source PDF page 14
\blockanchor{p014-b22}
\displayformula{a _ {n + 1} = \frac {a _ {n} + b _ {n}}{2}, b _ {n + 1} = \sqrt {b _ {n} a _ {n}} (n = 1, 2, \dots).}
% p014-b23; source PDF page 14
\blockanchor{p014-b23}
证明： $\lim_{n\to \infty}a_n$ 与 $\lim_{n\to \infty}b_n$ 存在且相等.\par
% p015-b00; source PDF page 15
\blockanchor{p015-b00}
\question{2}假设 $f(x)$ 在区间 $I$ 上二阶可导, 其中 $x_0 \in I$ 为 $f(x)$ 的极小值点, $f(x)$ 为 $I$ 上的严格凸函数, 那么 $f(x)$ 在 $I$ 上仅有一个唯一的极小值点.\par
% p015-b01; source PDF page 15
\blockanchor{p015-b01}
\question{3}证明: $\lim_{x \to +\infty} \frac{1}{x} \int_{0}^{x} f(t) \, \mathrm{d}t = \frac{1}{p} \int_{0}^{p} f(t) \, \mathrm{d}t$, 其中 $f(x)$ 是 $(- \infty, + \infty)$ 上的连续周期函数, 其中周期 $p > 0$.\par
% p015-b02; source PDF page 15
\blockanchor{p015-b02}
\question{4}求不定积分 $\int \sin^4 \mathrm{d}x$ 与 $\int \frac{x - 5}{x^2 - 3x + 4} \mathrm{d}x$.\par
% p015-b03; source PDF page 15
\blockanchor{p015-b03}
\question{5}设 $f(x, y) = |x - y| \varphi(x, y)$，其中 $\varphi(x, y)$ 在 (0, 0) 的某邻域内连续.\par
% p015-b04; source PDF page 15
\blockanchor{p015-b04}
(1) $\varphi(x, y)$ 满足什么条件时, $f_x(0, 0), f_y(0, 0)$ 存在?\par
% p015-b05; source PDF page 15
\blockanchor{p015-b05}
(2) $\varphi(x, y)$ 满足什么条件时, $f(x, y)$ 在 $(0, 0)$ 处可微?\par
% p015-b06; source PDF page 15
\blockanchor{p015-b06}
\question{6}求过点 $(2,1,3)$ 的平面, 使得其与 $x$ 轴, $y$ 轴, $z$ 轴围成的区域的体积为最小值.\par
% p015-b07; source PDF page 15
\blockanchor{p015-b07}
\question{7}求曲面 $z = \sqrt{3(x^2 + y^2)}$ 与平面 $x + y + z = 2$ 所围立体的表面积.\par
% p015-b08; source PDF page 15
\blockanchor{p015-b08}
\question{8}$S$ 为一个封闭曲面, $\pmb{n}$ 为 $S$ 的单位外法向量, $\pmb{r} = (x,y,z)$, 求证:\par
% p015-b09; source PDF page 15
\blockanchor{p015-b09}
\displayformula{\iint_{S}\frac{\cos(\widehat{\pmb{n},\pmb{r}})}{|\pmb{r}|^{2}}\mathrm{d}S=\left\{\begin{aligned}&0,&S\text{所围区域不包含原点,}\\ &4\pi,&S\text{所围区域包含原点.}\end{aligned}\right.}\par
% p015-b10; source PDF page 15
\blockanchor{p015-b10}
\question{9}设函数列 $\{f_n(x)\}$ 在 $[a, b]$ 上有连续的导函数, 且 $\{f_n(x)\}$ 在 $[a, b]$ 上收敛于 $f(x), \{f_n'(x)\}$ 在 $[a, b]$ 上一致收敛于 $g(x)$, 求证: $f(x)$ 在 $[a, b]$ 上处处可导, 且 $f'(x) = g(x)$.\par
% p015-b11; source PDF page 15
\blockanchor{p015-b11}
\question{10}求 $\sum_{n=1}^{\infty} \frac{n^2}{(1+a)^n}$, 其中 $a > 0$.\par
% p015-b12; source PDF page 15
\blockanchor{p015-b12}
\question{11}$\{f_n(x)\}$ 一致收敛于 $f(x)$, 且每个 $f_n(x)$ 都有界, 证明: $\{f_n(x)\}$ 在定义域中一致有界.\par
% p015-b13; source PDF page 15
\blockanchor{p015-b13}
\question{12}证明: $\sum_{n=1}^{\infty} ne^{-nx}$ 在 $(0, +\infty)$ 上不一致收敛, 但 $\sum_{n=1}^{\infty} ne^{-nx}$ 在 $(0, +\infty)$ 上连续.\par
